\chapter{Conclusions and Future Work}\label{conclusions}
% In this chapter you present all conclusions that can be drawn from the
% preceding chapters.
% It should not introduce new experiments, theories, investigations, etc.:
% These should have been written down earlier in the thesis.
% Therefore, conclusions can be brief and to the point.
\section{Conclusions}
In this thesis, we have formalised the call-by-box \lc, a refined modal \lc introduced by \cite{espritosanto_et_al:LIPIcs:2019:10525}.
We implemented the syntax, substitution, reduction and evaluation relations of \lab in Agda.
We also formalised Girard's and Gödel's embeddings which are embeddings from the standard call-by-name \lc (\lan) and call-by-value \lc (\lav) to the shared target, the call-by-box \lc (\lab).
For both embeddings, we have proven some properties that are preserved in \lab. We have proven the following propositions formally in Agda:

\begin{itemize}
  \item Girard's embedding preserves reduction (\cref{prop:girard-preservation-red}).
  \item Girard's embedding preserves evaluation (\cref{prop:girard-preservation-eval}).
  \item Gödel's embedding preserves reduction (\cref{prop:godel-preservation-red} and \ref{prop:godel-preservation-red2s}).
\end{itemize}

To prove these propositions, we relied on several helper lemmas.
The most important are \cref{lem:girard-pres-subst} for the properties about Girard's translation,
and \cref{lem:godel-pres-subst} for the properties about Gödel's translation.
These helper lemmas state that substitution is preserved and reflected by the translation.
They were the most complex to prove, but they are fundamental for the proofs about preservation of reduction and evaluation.

We also saw that the formal definition of \raiseembn in Gödel's embedding is more complex than the less formal definition due to the use of the rename function.
This rename function made proofs about Gödel's embedding significantly more complex than those about Girard's embedding,
which is one reason why we did not prove the preservation of evaluation for Gödel's embedding formally. 

In conclusion, this thesis provides a formal implementation for \lab in Agda. 
It contains the syntax, substitution, reduction and evaluation relations as well as some proofs of Girard's and Gödel's embeddings.
The formalisation could serve as a basis for further research into the unification of call-by-name and call-by-value evaluation strategies in a modal setting,
and could also serve as a foundation for future programming languages or proof assistants that natively support both evaluation strategies.

\section{Future Work}
This thesis has formalised the core functionalities of \lab and has proven some important properties of Girard's and Gödel's embeddings.
However, there are still many areas for future work regarding this formalisation. We list some opportunities for future work below.

The formalisation of \lab can be extended with more features, such as case expressions, natural numbers or recursion.
The book PLFA \cite{plfa22.08} implements these features in the standard \lc, and they can be implemented in \lab as well.

Preservation of evaluation under Gödel's embedding has not been proven formally yet, but it is an interesting property to prove.
It would make the formalisation more complete and could be useful for future research.

We can extend the formalisation with more properties, such as reflection of reduction and evaluation under Girard's and Gödel's embeddings.
This would show that the target \lab, does not simulate any other reductions or evaluations than those of \lan and \lav.
For example, the reflection of reduction under Girard's translation can be stated as follows: if $\girard{M} \redbbox \girard{N}$, then $M \rednb N$ in \lan.

We only formalised \lab in Agda, but it could be interesting to also formalise different frameworks that can simulate both \textsf{cbn} and \textsf{cbv},
such as the linear logic approach by Ehrhard and Giulio Guerrieri \cite{EhrhardGuerrieri2016} or Maraist et al. \cite{MARAIST1999175}, 
and the thunk-based approach by Hatcliff and Danvy \cite{HATCLIFF_DANVY_1997}.
All these formalisations could then be compared to see which is simpler, more efficient or more powerful.

The boxed \lc can also be used to construct an interpreter for a programming language that supports both \textsf{cbn} and \textsf{cbv} evaluation.
This is a more practical application and it could be interesting to see \lab in practice.  